\chapter{通信逻辑}

通信逻辑是 CCB 的核心。用户看到的是 \cmd{ccb ask agent1 ...}、\cmd{ccb watch job_xxx}、\cmd{ccb inbox talk1}、\cmd{ccb ack talk1}，源码里实际流动的是 envelope、job、message、attempt、inbound event、reply、callback edge 和 completion decision。

\section{总览}

一次普通 ask 的端到端路径如下。

\begin{Verbatim}[fontsize=\small,frame=single]
CLI tokens
  -> parse_ask
  -> submit_ask
  -> MessageEnvelope
  -> ccbd submit handler
  -> dispatcher submission plan
  -> JobRecord / JobEvent / dispatcher queue
  -> MessageRecord / AttemptRecord / InboundEventRecord
  -> dispatcher tick starts running job
  -> provider execution service
  -> completion polling and tracker
  -> finalization
  -> AttemptRecord / ReplyRecord / reply inbound event
  -> watch / queue / inbox / trace views
\end{Verbatim}

这条链路有三个关键分界。

\begin{enumerate}
  \item \src{MessageEnvelope} 是 CLI 到 daemon 的边界对象。
  \item \src{JobRecord} 是 dispatcher 执行事实。
  \item \src{MessageRecord}、\src{AttemptRecord}、\src{ReplyRecord} 和 \src{InboundEventRecord} 是通信事实。
\end{enumerate}

调试时应先判断问题落在哪个边界：CLI 没提交、daemon 没接收、dispatcher 没排队、provider 没执行、completion 没终止、reply 没记录、mailbox 没投影，还是用户视图没渲染。

\section{CLI ask 解析}

ask 解析入口在 \src{lib/cli/parser_runtime/ask.py}，提交服务在 \src{lib/cli/services/ask.py} 和 \src{lib/cli/services/ask_runtime/submission.py}。它们处理以下职责：

\begin{itemize}
  \item 解析 \cmd{--compact}、\cmd{--silence}、\cmd{--callback}、\cmd{--artifact-request}、\cmd{--artifact-reply}、\cmd{--artifact-io}。
  \item 解析 target 和 message body，支持 stdin 追加。
  \item 推断 sender。当前 workspace agent 能识别时使用 agent，否则 fallback 到 user。
  \item 校验 target 和 sender 是否存在于项目配置。
  \item 在普通 ask 中追加 reply guidance，除非用户已经给出明确输出要求。
  \item 在大 request 或显式 artifact flag 下把 request body 写入 text artifact。
  \item 构造 \src{MessageEnvelope} 并通过 mounted daemon client 提交。
\end{itemize}

ask flags 不应被理解为显示格式开关。它们是路由和交付策略。

\begin{longtable}{p{0.25\linewidth}p{0.65\linewidth}}
\toprule
flag & 语义 \\
\midrule
\cmd{--compact} & 请求子 agent 返回提炼结果，同时保留关键信息。 \\
\cmd{--silence} & 成功时静默，失败或阻塞仍应暴露。适合无需结果的独立工作。 \\
\cmd{--callback} & 把子任务结果作为 continuation 送回当前 agent。适合 parent 需要 child 结果才能完成的任务。 \\
\cmd{--artifact-request} & 强制 request body 存成 CCB text artifact。 \\
\cmd{--artifact-reply} & 强制 final reply 存成 CCB text artifact。 \\
\cmd{--artifact-io} & 同时启用 request 和 reply artifact。 \\
\bottomrule
\end{longtable}

\section{MessageEnvelope}

\src{lib/ccbd/api_models_runtime/messages.py} 中的 \src{MessageEnvelope} 负责规范化 daemon API 边界。submit handler 位于 \src{lib/ccbd/handlers/submit.py}，它从 RPC payload 重建 envelope，然后调用 \src{dispatcher.submit(envelope)}。

envelope 保留：

\begin{itemize}
  \item \src{project_id}、\src{to_agent}、\src{from_actor}；
  \item \src{body} 和可选 \src{body_artifact}；
  \item \src{message_type}、\src{delivery_scope}；
  \item \src{silence_on_success} 和 \src{route_options}；
  \item \src{task_id}、\src{reply_to}。
\end{itemize}

delivery scope 的验证在 API model 层存在，不只依赖 CLI。这可以防止其他 client 绕过 CLI 后提交结构不一致的 payload。

\section{dispatcher submission}

dispatcher facade 位于 \src{lib/ccbd/services/dispatcher_runtime/facade.py}。提交路径由 \src{submission_service.py}、\src{submission_recording.py} 和 \src{lifecycle.py} 共同实现。

提交时，dispatcher 先生成 submission plan：

\begin{enumerate}
  \item 验证 sender。
  \item 验证 body artifact。
  \item 验证 callback request。
  \item 解析 target agents。
  \item 确认 target availability。
  \item 为每个 target 生成 \src{_JobDraft}。
\end{enumerate}

然后 \src{_submit_plan} 创建 job id、写入 \src{JobRecord}、写入 \src{job_accepted} 或 \src{job_queued} event、更新 dispatcher state，并调用 message bureau 记录 submission。

广播 ask 会产生一个 submission id 和多个 jobs。单 target ask 不需要把用户认知绑定到 submission id，但内部仍可通过 message-bureau 记录追踪 lineage。

\section{message bureau 记录}

message bureau 的 submission recording 负责把 dispatcher job 映射成通信记录：

\begin{itemize}
  \item 一个 \src{MessageRecord} 表示逻辑请求。
  \item 每个 job 对应一个 \src{AttemptRecord}。
  \item 每个 target agent 收到一个 \src{InboundEventRecord}，事件类型通常为 \src{task_request}。
\end{itemize}

当 dispatcher 启动 job 时，\src{mark_attempt_started} 会把 attempt 移到 running，claim mailbox inbound event，并把 message state 更新到 running。

当 job 进入 terminal 时，message bureau 更新 attempt state、消费或 abandon inbound event、刷新 message state，并在需要时写入 reply record 和 caller mailbox reply event。

\section{运行、轮询和 finalization}

运行路径从 \src{tick_jobs} 开始。dispatcher 对每个可运行 slot 启动下一个 queued job。核心函数分布在：

\begin{itemize}
  \item \src{lib/ccbd/services/dispatcher_runtime/lifecycle_start_runtime/tick.py}
  \item \src{lib/ccbd/services/dispatcher_runtime/lifecycle_start_runtime/start.py}
  \item \src{lib/ccbd/services/dispatcher_runtime/polling_service.py}
  \item \src{lib/ccbd/services/dispatcher_runtime/finalization_runtime/service.py}
  \item \src{lib/ccbd/services/dispatcher_runtime/finalization_runtime/message_bureau.py}
\end{itemize}

\src{start_running_job} 持久化 running state，标记 message-bureau attempt started，写 completion snapshot，更新 dispatcher active state，同步 runtime busy state，并在 runtime binding 可执行时启动 provider execution。

\src{poll_completion_updates} 读取 provider execution update，追加 completion item event，ack execution item，更新 completion tracker。如果 tracker 给出 terminal decision，dispatcher 进入 \src{complete_job}。

finalization 做的事比“保存 reply”更多：

\begin{itemize}
  \item 持久化 terminal completion。
  \item 记录 message-bureau terminal attempt。
  \item 判断自动 retry。
  \item 处理 reply artifact spill。
  \item 写入 reply record。
  \item 解析 callback child continuation。
  \item 解析 reply delivery terminal state。
\end{itemize}

\section{watch}

\cmd{watch} 是轮询型观察命令。CLI 服务 \src{lib/cli/services/watch.py} 读取环境变量：

\begin{itemize}
  \item \src{CCB_WATCH_TIMEOUT_S}，默认 10 秒。
  \item \src{CCB_WATCH_POLL_INTERVAL_S}，默认 0.1 秒。
\end{itemize}

\src{lib/cli/services/watch_runtime.py} 使用 cursor 调用 \src{client.watch(target, cursor=cursor)}。每次 daemon 返回一个 batch，CLI yield 非空 events，并更新 cursor。如果 batch terminal，则退出。

watch 的重要可靠性机制是 persisted terminal fallback。当 daemon 连接失败或超时时，CLI 会尝试从持久化终态 watch payload 恢复最后结果。这样可以缓解 daemon 重启窗口导致的用户可见中断。

\section{queue 和 inbox}

\cmd{queue} 和 \cmd{inbox} 不是直接读取 dispatcher state。它们走 message-bureau control queue 层：

\begin{itemize}
  \item \src{lib/message_bureau/control_queue_runtime/views_runtime/summary.py}
  \item \src{lib/message_bureau/control_queue_runtime/views_runtime/agent.py}
  \item \src{lib/message_bureau/control_queue_runtime/views_runtime/inbox.py}
  \item \src{lib/message_bureau/control_queue_runtime/events.py}
\end{itemize}

\cmd{queue all} 聚合每个 agent 的 mailbox summary，展示 agent count、queued agent count、active agent count、total queue depth 和 total pending reply count。dispatcher facade 会额外叠加 runtime state 和 runtime health。

\cmd{queue <agent>} 展示单个 agent 的 mailbox state、lease version、queue depth、pending reply count、active inbound event 和可选 queued events。如果 mailbox summary 缺失或损坏，视图会 degraded，但 detail 仍可以从 inbound event records 重建一部分信息。

\cmd{inbox <agent>} 展示 mailbox head。head 如果是 reply，会被补充 reply id、source actor、terminal status、notice、job id、finished time、progress time 和 reply text。detail 模式列出 pending items。

\section{ack}

\cmd{ack <agent> [inbound_event_id]} 的语义非常严格。实现位于 \src{lib/message_bureau/control_queue_runtime/ack.py}。

ack 只允许确认 head event。它支持两类 head：

\begin{enumerate}
  \item \src{task_reply}，即 caller inbox 收到的 reply。
  \item 已经 terminal 的 \src{task_request}，通常用于清理已经取消或终止但仍在 head 的请求事件。
\end{enumerate}

如果用户指定的 inbound event 不是 head，ack 会拒绝并提示当前 head。这个限制是故意的，它保证 inbox 处理顺序不会被跳过。

ack 还会拒绝已经安排 automatic reply delivery 的 reply，因为这类 reply 后续由 delivery job 接管。手动 ack 可能破坏 delivery head 的一致性。

\section{trace}

\cmd{trace} 是通信 lineage 的标准入口。实现位于 \src{lib/message_bureau/control_trace_runtime/service.py}。它根据 id 前缀解析目标类型：

\begin{itemize}
  \item \src{sub_}：submission。
  \item \src{msg_}：message。
  \item \src{att_}：attempt。
  \item \src{rep_}：reply。
  \item \src{job_}：job。
\end{itemize}

trace 返回统一 payload，包含 selected object、counts 和关联数组：messages、attempts、replies、events、jobs。submission trace 会收集该 submission 的最新 messages，再收集 attempts、replies、events 和 job summaries。job trace 会先通过 job id 找 attempt，再找 message 和 replies。

这说明 trace 的图模型是：

\begin{Verbatim}[fontsize=\small,frame=single]
submission
  -> message
    -> attempt
      -> job
      -> inbound events
      -> reply
\end{Verbatim}

用户手册应建议：只要手里有任意 \src{job_id}、\src{message_id} 或 \src{reply_id}，先用 \cmd{trace} 重建上下文，再决定 retry、resubmit、ack 还是 doctor。

\section{callback}

callback 不是同步等待。callback mode 由 \src{route_options["mode"] == "callback"} 表示，核心实现位于 \src{lib/ccbd/services/dispatcher_runtime/callbacks.py} 和 \src{lib/message_bureau/callback_edges.py}。

active CCB task 内发起普通 nested ask 会被拒绝，除非它使用 \cmd{--callback} 或 \cmd{--silence}。这是因为 parent agent 如果需要 child 结果，应当显式建立 callback edge；如果不需要结果，应当显式选择 fire-and-forget。

callback 的限制包括：

\begin{itemize}
  \item 必须 exactly one target。
  \item 必须存在 active parent job。
  \item 必须存在 parent message。
  \item message-bureau support 必须可用。
  \item parent 不能已有 outstanding callback。
  \item callback continuation job 不能再 \cmd{--callback} 回该 continuation 的 original caller；它必须直接完成当前 turn，由 CCB 把 completion 继续向上游投递。
\end{itemize}

child 完成后，CCB 不把 child reply 直接投递给 original caller。它记录 child reply，向 parent agent 提交 continuation job，并更新 callback edge。continuation job 的 final answer 才是上游投递路径；agent 不应再用 ask 把 final result 发送给 original caller。repair 和 timeout sweep 可以恢复或失败 pending callback edges。

\section{retry 和 resubmit}

\cmd{retry} 和 \cmd{resubmit} 都让任务再次执行，但语义不同。

\cmd{resubmit <message_id>} 重新规划一条 message resubmission，走普通 submit plan，返回 original message id、新 message id、submission id 和新 jobs。它创建新的 message/submission chain。

\cmd{retry <job_id|attempt_id>} 保留原 message lineage。它要求目标 attempt 已 terminal，拒绝 completed attempt，并要求该 attempt 是对应 agent 的最新 attempt。然后它创建新 job，记录 retry attempt，并把新 job 追加到原 submission。

ask retry 有一个细节：如果旧 job 已经看到 provider anchor 或已经开始 reply，新 retry 可能把 request body 改成 \src{continue}，并通过 provider options 标记 retry delivery mode。这样可以尽量利用 provider resume 语义，而不是把同一 prompt 完全重复投递。

自动 retry 与手动 retry 分离。自动 retry 依赖 message retry policy、失败 reason/error type、provider resume 支持和 max attempts。实现入口在 \src{lib/ccbd/services/dispatcher_runtime/finalization_retry_runtime/}。

\section{cancel}

\cmd{cancel <job_id>} 由 \src{lib/ccbd/services/dispatcher_runtime/cancellation.py} 实现。取消流程包括：

\begin{enumerate}
  \item 校验 job 存在且尚未 terminal。
  \item 写入 cancel requested。
  \item 从 dispatcher queued/active state 移除。
  \item 调用 execution service cancel。
  \item 构造 cancelled completion decision。
  \item 写 completion snapshot 和 terminal event。
  \item 记录 message-bureau terminal state。
  \item 解析 reply-delivery terminal state。
\end{enumerate}

取消不是简单从队列删除。它必须把执行层和通信层都推进到一致终态，否则 trace、inbox、queue 和 watch 会互相矛盾。

\section{artifact 和 reply delivery}

terminal reply 可能通过 \src{lib/ccbd/services/dispatcher_runtime/finalization_runtime/artifacts.py} 写入 text artifact。触发条件是用户显式指定 \cmd{--artifact-reply}，或 reply 大于 4 KiB。decision diagnostics 会记录 \src{reply_artifact}，visible reply 变成 artifact stub。

reply delivery 是独立 subsystem，入口在 \src{lib/ccbd/services/dispatcher_runtime/reply_delivery.py} 和 \src{reply_delivery_runtime/}。它负责准备、claim、start completion、terminal、head 和 repair。它解释了为什么 ack 不能随意消费已经安排 automatic reply delivery 的 reply head。

\section{通信调试顺序}

开发者排查通信问题时可以按以下顺序。

\begin{enumerate}
  \item 用 \cmd{trace <job_id>} 确认 job 是否绑定 message/attempt/reply。
  \item 用 \cmd{queue --detail <agent>} 看 target agent mailbox 是否仍有 task request。
  \item 用 \cmd{inbox --detail <caller>} 看 caller 是否已有 reply head。
  \item 用 \cmd{watch <job_id>} 看 completion event 是否 terminal。
  \item 用 \cmd{doctor storage} 判断 message-bureau/mailbox summary 是否 degraded。
  \item 如果 attempt failed 且是最新终态，选择 \cmd{retry}；如果需要重新发整条请求，选择 \cmd{resubmit}。
  \item 如果 head reply 已读且未被 automatic delivery 接管，用 \cmd{ack} 消费。
\end{enumerate}

这套顺序的原则是先重建 lineage，再处理队列，最后改变状态。直接 cancel、retry 或 ack 可能掩盖真正的断点。
